\documentclass[11pt]{article}

\usepackage[T1]{fontenc}
\usepackage{amsmath,amssymb,amsthm}
\usepackage[margin=1in]{geometry}
\usepackage[hidelinks]{hyperref}
\setlength{\emergencystretch}{3em}

\newtheorem{theorem}{Theorem}
\newtheorem{lemma}[theorem]{Lemma}
\newtheorem{corollary}[theorem]{Corollary}
\newtheorem{proposition}[theorem]{Proposition}
\newtheorem{conjecture}[theorem]{Conjecture}
\theoremstyle{remark}
\newtheorem{remark}[theorem]{Remark}

\newcommand{\R}{\mathbb R}
\newcommand{\PP}{\mathbb P}
\newcommand{\Per}{\operatorname{Per}}

\title{The Critical Weak-\(L^3\) Clock Forces an Ultraviolet Increment}
\author{John Robert Lawson and Codex}
\date{Research round ending 25 July 2026}

\begin{document}
\maketitle

\begin{abstract}
This round closes the low-frequency side of the spectral dichotomy for
an energy-efficient exact \(q=4\) first-record branch.  Finite-volume
weak-\(L^3\) capacity makes the velocity on a future selected amplitude
layer almost orthogonal to a fixed earlier record.  A new projected
low-pass estimate,
\[
 \|P_{\leq\Lambda}(u(t)-u(s))\|_2
 \lesssim
 \bigl(M^2\Lambda^{3/2}
       +\nu\sqrt E\,\Lambda^2\bigr)(t-s),
\]
then shows that the resulting fixed signed change cannot occur below
\[
 \Lambda_j\asymp
 m_j^{7/3}j^{2/3}
 \asymp R_j^{-7/6}j^{2/3}.
\]
A fixed \(L^2\) ultraviolet floor follows at one endpoint, together with
\[
 \|\nabla u(s_j)\|_2^2
 \gtrsim R_j^{-7/3}j^{4/3}.
\]
For any Leray--Hopf continuation through the first singular time, this
forces a fixed downward kinetic-energy jump.  Hence energy equality at
that time conditionally excludes every energy-efficient exact \(q=4\)
branch.  General Leray--Hopf energy equality remains open, so this is a
conditional reduction rather than a Clay resolution.  Static packet and
driver families show sharpness of the spatial and instantaneous clock
powers under the inputs used; all presently available record-gap weights
remain summable.
\end{abstract}

\section{Epistemic scope}

All statements about the selected layer are conditional on one putative
smooth finite-energy Navier--Stokes trajectory satisfying the exact
\(q=4\) first-record ledger.  The low-pass lemma, ultraviolet increment,
terminal energy-defect corollary, exponent ledger, and sharpness examples
are repository deductions and await external mathematical review.

The statement that energy equality for arbitrary Leray--Hopf solutions
remains open, as well as representative sufficient Besov conditions, is
recorded by Cheskidov and Luo~\cite{CL20}.  No such sufficient condition
is derived from the hypotheses below.  The packet families in
Section~\ref{sec:sharpness} vary with scale and are not one
Navier--Stokes evolution.  The Clay problem remains unsolved.

\section{Exact record setting and Fourier split}

Let \(u\) be a smooth divergence-free solution on
\(\R^3\times[0,T^*)\) of
\[
 \partial_tu+u\cdot\nabla u=-\nabla p+\nu\Delta u,
 \qquad
 \nabla\cdot u=0,
\]
where \(T^*<\infty\) is its first possible singular time.  Assume
\begin{equation}
 \label{eq:energy-bound}
 \sup_{t<T^*}\|u(t)\|_2^2\leq E_0.
\end{equation}

Let \(t_j\uparrow T^*\) be exact \(q=4\) first-record times for
\[
 m_j:=\|u(t_j)\|_{L^{3,\infty}}.
\]
Thus
\begin{equation}
 \label{eq:q4}
 m_j\asymp2^{2j},
 \qquad
 t_{j+1}-t_j\asymp\frac{2^{-11j}}j,
 \qquad
 \|u(t)\|_{L^{3,\infty}}\leq m_j
 \quad(0\leq t\leq t_j).
\end{equation}
At \(t_j\), choose the near-extremising amplitude layer
\[
 A_j:=\{a_j<|u(t_j)|\leq2a_j\},
 \qquad
 e_j:=\int_{A_j}|u(t_j)|^2\,dx,
 \qquad
 R_j:=|A_j|^{1/3}.
\]
The energy-efficient branch is
\begin{equation}
 \label{eq:layer-energy}
 0<c_0\leq e_j\leq E_0,
 \qquad
 R_j\asymp m_j^{-2}.
\end{equation}

Choose a real radial \(\psi\in C_c^\infty(\R^3)\) such that
\[
 0\leq\psi\leq1,\qquad
 \psi=1\quad\hbox{on }B_1,\qquad
 \psi=0\quad\hbox{outside }B_2,
\]
and define
\begin{equation}
 \label{eq:LP}
 \widehat{P_{\leq\Lambda}f}(\xi)
 =\psi(\xi/\Lambda)\widehat f(\xi),
 \qquad
 P_{>\Lambda}:=I-P_{\leq\Lambda}.
\end{equation}
Both are self-adjoint \(L^2\) contractions.  Since the second multiplier
vanishes for \(|\xi|\leq\Lambda\),
\begin{equation}
 \label{eq:bernstein-high}
 \|P_{>\Lambda}f\|_2
 \leq\Lambda^{-1}\|\nabla f\|_2.
\end{equation}

\section{The two endpoint lemmas}

\begin{lemma}[Finite-volume weak-\(L^3\) capacity]
\label{lem:capacity}
For every measurable \(E\subset\R^3\) of finite measure and every
\(g\in L^{3,\infty}(\R^3)\),
\begin{equation}
 \label{eq:capacity}
 \int_E|g|^2\,dx
 \leq C|E|^{1/3}\|g\|_{L^{3,\infty}}^2.
\end{equation}
\end{lemma}

\begin{proof}
Put \(M=\|g\|_{L^{3,\infty}}\) and
\(\lambda_0=M|E|^{-1/3}\).  The distribution bound
\[
 |\{|g|>\lambda\}|\leq M^3\lambda^{-3}
\]
and the layer-cake formula give
\begin{align*}
 \int_E|g|^2\,dx
 &=
 \int_0^\infty
 2\lambda\,|E\cap\{|g|>\lambda\}|\,d\lambda\\
 &\leq
 \lambda_0^2|E|
 +2M^3\int_{\lambda_0}^{\infty}\lambda^{-2}\,d\lambda\\
 &\leq3M^2|E|^{1/3}.
\end{align*}
\end{proof}

\begin{lemma}[Critical low-pass clock]
\label{lem:clock}
Let \(0\leq s<t<T^*\), and suppose
\[
 \sup_{s\leq\tau\leq t}
 \|u(\tau)\|_{L^{3,\infty}}\leq M.
\]
Then
\begin{equation}
 \label{eq:clock}
 \boxed{
 \|P_{\leq\Lambda}(u(t)-u(s))\|_2
 \leq
 C\bigl(
 M^2\Lambda^{3/2}
 +\nu\sqrt{E_0}\Lambda^2
 \bigr)(t-s).
 }
\end{equation}
\end{lemma}

\begin{proof}
The projected equation is
\begin{equation}
 \label{eq:projected}
 \partial_tu
 =-\PP\nabla\cdot(u\otimes u)+\nu\Delta u.
\end{equation}
For a tensor \(F\in L^{3/2,\infty}\), Lorentz duality gives
\begin{equation}
 \label{eq:dual}
 \|P_{\leq\Lambda}\PP\nabla\cdot F\|_2
 \leq
 C\|F\|_{3/2,\infty}
 \sup_{\|h\|_2=1}
 \|\nabla\PP P_{\leq\Lambda}h\|_{3,1}.
\end{equation}
Plancherel and Fourier Cauchy--Schwarz imply
\[
 \|\nabla\PP P_{\leq\Lambda}h\|_2
 \leq C\Lambda\|h\|_2,
\qquad
 \|\nabla\PP P_{\leq\Lambda}h\|_\infty
 \leq C\Lambda^{5/2}\|h\|_2.
\]
Real interpolation between these two estimates yields
\begin{equation}
 \label{eq:lorentz-bernstein}
 \|\nabla\PP P_{\leq\Lambda}h\|_{3,1}
 \leq C\Lambda^{3/2}\|h\|_2.
\end{equation}
Since
\[
 \|u\otimes u\|_{3/2,\infty}
 \leq C\|u\|_{3,\infty}^2,
\]
the nonlinear term in \eqref{eq:projected}, after low-pass projection,
has \(L^2\) norm at most \(CM^2\Lambda^{3/2}\).  The viscous term obeys
\[
 \|P_{\leq\Lambda}\Delta u\|_2
 \leq C\Lambda^2\|u\|_2
 \leq C\sqrt{E_0}\Lambda^2.
\]
Integrating \eqref{eq:projected} from \(s\) to \(t\) proves
\eqref{eq:clock}.
\end{proof}

\begin{remark}
The energy-only compact-carrier clock has a
\(\Lambda^{5/2}\) nonlinear rate.  Lemma~\ref{lem:clock} gains one
frequency power by using the first-record weak-\(L^3\) ceiling.  The
estimate controls a global \(L^2\) increment; no spatial carrier
compactness is assumed.
\end{remark}

\section{The fixed ultraviolet increment}

\begin{theorem}[Critical-clock ultraviolet increment]
\label{thm:increment}
Under \eqref{eq:energy-bound}--\eqref{eq:layer-energy}, there are a fixed
integer \(L\geq1\), a fixed \(\kappa>0\), and constants \(c>0\) and
\(j_0\) such that, with
\begin{equation}
 \label{eq:block-gap}
 \Delta_j^{(L)}:=t_j-t_{j-L}
 \asymp_L\frac{m_j^{-11/2}}j
\end{equation}
and
\begin{equation}
 \label{eq:Lambda}
 \Lambda_j
 :=\kappa\bigl(m_j^2\Delta_j^{(L)}\bigr)^{-2/3},
\end{equation}
one has, for every \(j\geq j_0\),
\begin{equation}
 \label{eq:signed-high}
 \boxed{
 \left\langle
 P_{>\Lambda_j}\bigl(u(t_j)-u(t_{j-L})\bigr),
 u(t_j)\mathbf 1_{A_j}
 \right\rangle
 \geq c,
 }
\end{equation}
and
\begin{equation}
 \label{eq:norm-high}
 \boxed{
 \left\|
 P_{>\Lambda_j}\bigl(u(t_j)-u(t_{j-L})\bigr)
 \right\|_2
 \geq c.
 }
\end{equation}
\end{theorem}

\begin{proof}
Set \(i=j-L\) and
\[
 f_j:=u(t_j)\mathbf 1_{A_j}.
\]
Then \(\|f_j\|_2^2=e_j\).  Lemma~\ref{lem:capacity},
\(m_{j-L}/m_j\lesssim2^{-2L}\), and \(R_jm_j^2\asymp1\) give
\begin{equation}
 \label{eq:early-empty}
 \int_{A_j}|u(t_i)|^2\,dx
 \leq CR_jm_i^2
 \leq C2^{-4L}.
\end{equation}
Choose \(L\) sufficiently large, depending only on the fixed ledger
constants, so that Cauchy--Schwarz implies
\begin{equation}
 \label{eq:old-pair}
 |\langle u(t_i),f_j\rangle|
 \leq\frac{c_0}{4}.
\end{equation}
On the other hand,
\[
 \langle u(t_j),f_j\rangle=e_j\geq c_0.
\]
Consequently,
\begin{equation}
 \label{eq:total-pair}
 \langle u(t_j)-u(t_i),f_j\rangle
 \geq\frac{3c_0}{4}
 =:c_1.
\end{equation}

The first two relations in \eqref{eq:q4} give
\eqref{eq:block-gap} by summing a fixed number of geometrically
separated gaps.  The first-record property gives
\[
 \sup_{t_i\leq t\leq t_j}
 \|u(t)\|_{3,\infty}\leq m_j.
\]
Apply Lemma~\ref{lem:clock} on this interval.  Definition
\eqref{eq:Lambda} gives the exact nonlinear identity
\begin{equation}
 \label{eq:nonlinear-small}
 m_j^2\Lambda_j^{3/2}\Delta_j^{(L)}
 =\kappa^{3/2},
\end{equation}
whereas
\begin{equation}
 \label{eq:viscous-small}
 \Delta_j^{(L)}\Lambda_j^2
 \asymp_Lm_j^{-5/6}j^{1/3}
 \longrightarrow0.
\end{equation}
Choose \(\kappa\) sufficiently small and then \(j\) sufficiently large.
Equations \eqref{eq:clock}, \eqref{eq:nonlinear-small}, and
\eqref{eq:viscous-small} then yield
\begin{equation}
 \label{eq:low-small}
 \|P_{\leq\Lambda_j}(u(t_j)-u(t_i))\|_2
 \leq\frac{c_1}{2\sqrt{E_0}}.
\end{equation}
Pair \eqref{eq:low-small} with \(f_j\) and subtract the result from
\eqref{eq:total-pair}.  This proves \eqref{eq:signed-high}.
Cauchy--Schwarz, together with \(\|f_j\|_2\leq\sqrt{E_0}\), proves
\eqref{eq:norm-high}.
\end{proof}

\begin{corollary}[Endpoint energy and enstrophy floors]
\label{cor:endpoint}
For each sufficiently large \(j\), at least one
\[
 s_j\in\{t_{j-L},t_j\}
\]
satisfies
\begin{equation}
 \label{eq:endpoint-high}
 \|P_{>\Lambda_j}u(s_j)\|_2\geq c
\end{equation}
and
\begin{equation}
 \label{eq:endpoint-enstrophy}
 \boxed{
 \|\nabla u(s_j)\|_2^2
 \gtrsim
 \Lambda_j^2
 \asymp_L
 m_j^{14/3}j^{4/3}
 \asymp_L
 R_j^{-7/3}j^{4/3}.
 }
\end{equation}
The associated physical length is
\begin{equation}
 \label{eq:rclock}
 r_j:=\Lambda_j^{-1}
 \asymp_L
 R_j^{7/6}j^{-2/3}.
\end{equation}
\end{corollary}

\begin{proof}
The triangle inequality applied to \eqref{eq:norm-high} gives
\eqref{eq:endpoint-high} at one endpoint.  Equation
\eqref{eq:bernstein-high} then gives the first inequality in
\eqref{eq:endpoint-enstrophy}.  Finally,
\[
 \Lambda_j
 \asymp_Lm_j^{7/3}j^{2/3}
 \asymp_LR_j^{-7/6}j^{2/3}
\]
follows from \eqref{eq:block-gap}, \eqref{eq:Lambda}, and
\(R_j\asymp m_j^{-2}\).
\end{proof}

\begin{remark}
The scale \(r_j\) lies strictly below the previous sharp Morrey
fragmentation scale \(R_j^{27/26+o(1)}\).  The conclusion is a fixed
high-pass \(L^2\) floor, not merely a large first derivative.
\end{remark}

\section{The forced terminal energy defect}

\begin{corollary}[Anomalous loss at the first singular time]
\label{cor:defect}
Suppose the smooth solution has a Leray--Hopf continuation through
\(T^*\), represented weakly continuously in \(L^2\), and put
\[
 u_*:=u(T^*).
\]
Then
\begin{equation}
 \label{eq:defect}
 \boxed{
 \lim_{t\uparrow T^*}\|u(t)\|_2^2-\|u_*\|_2^2
 \geq c>0.
 }
\end{equation}
Consequently, energy equality at \(T^*\), equivalently strong left
\(L^2\) continuity there, excludes the energy-efficient exact \(q=4\)
branch.
\end{corollary}

\begin{proof}
Both possible endpoint sequences converge to \(T^*\).  Pass to an
infinite subsequence using the same endpoint choice.  Weak continuity
gives
\begin{equation}
 \label{eq:weak}
 u(s_j)\rightharpoonup u_*
 \quad\hbox{in }L^2.
\end{equation}
Since \(\Lambda_j\to\infty\) and \(u_*\in L^2\),
\begin{equation}
 \label{eq:trace-tail}
 \|P_{>\Lambda_j}u_*\|_2\longrightarrow0.
\end{equation}
Equations \eqref{eq:endpoint-high} and \eqref{eq:trace-tail} give
\begin{equation}
 \label{eq:strong-defect}
 \|u(s_j)-u_*\|_2
 \geq
 \|P_{>\Lambda_j}(u(s_j)-u_*)\|_2
 \geq c-o(1).
\end{equation}

The smooth energy equality before \(T^*\) makes
\(\|u(t)\|_2^2\) nonincreasing, so its left limit \(E_-\) exists.
Weak convergence gives
\[
 \|u(s_j)-u_*\|_2^2
 =
 \|u(s_j)\|_2^2-\|u_*\|_2^2+o(1).
\]
Taking the limit in \eqref{eq:strong-defect} proves
\(E_--\|u_*\|_2^2\geq c\).

If energy equality held through \(T^*\), taking the limit in the smooth
preterminal identity would instead give
\[
 E_-+2\nu\int_0^{T^*}\|\nabla u\|_2^2\,dt
 =
 \|u_*\|_2^2
 +2\nu\int_0^{T^*}\|\nabla u\|_2^2\,dt,
\]
which contradicts \eqref{eq:defect}.
\end{proof}

\begin{remark}
Corollary~\ref{cor:defect} does not assert an unconditional
contradiction.  It identifies the remaining branch with the unresolved
possibility of anomalous energy loss by a Leray--Hopf solution.  The
general energy-equality question is explicitly open in
\cite{CL20}; that paper proves equality under additional weak-in-time
Besov hypotheses.  No claim that the present record schedule meets those
hypotheses is made.
\end{remark}

\section{Sharp static ledgers}
\label{sec:sharpness}

\begin{proposition}[Packet-cloud sharpness]
\label{prop:packets}
Let \(m\to\infty\), retain the corresponding index \(j\), and set
\begin{equation}
 \label{eq:packet-ledger}
 R=m^{-2},
 \qquad
 r=m^{-7/3}j^{-2/3},
 \qquad
 a=m^3,
 \qquad
 N\asymp(R/r)^3\asymp mj^2.
\end{equation}
There are smooth compactly supported divergence-free fields \(U_j\)
for which
\begin{equation}
 \label{eq:packet-norms}
 \|U_j\|_2^2\asymp1,
 \qquad
 \|U_j\|_{3,\infty}\asymp m,
\end{equation}
a fixed regular amplitude band has volume \(\asymp R^3\) and energy
\(\asymp1\), and
\begin{equation}
 \label{eq:packet-costs}
 \|\nabla U_j\|_2^2
 \asymp r^{-2}
 \asymp m^{14/3}j^{4/3},
 \qquad
 \Per(A_j)\asymp R^{11/6}j^{2/3}.
\end{equation}
For every fixed \(0<\beta<1/9\),
\begin{equation}
 \label{eq:packet-morrey}
 \sup_x\int_{B_\rho(x)}|U_j|^2\,dx
 \leq C_\beta\rho^\beta
 \qquad(0<\rho\leq1).
\end{equation}
\end{proposition}

\begin{proof}
Choose a nonzero
\(\phi\in C_c^\infty(B_1;\R^3)\) with
\(\nabla\cdot\phi=0\) and a regular amplitude band of positive measure.
Place \(N\) centres quasi-uniformly in one fixed cube, with separation
\(d\asymp N^{-1/3}\), and define
\[
 U_j(x):=
 a\sum_{k=1}^N
 \phi\left(\frac{x-x_k}{r}\right).
\]
Since \(r/d\to0\), the supports are disjoint.  Direct scaling gives
\[
 \|U_j\|_2^2\asymp Na^2r^3\asymp1,
 \qquad
 \|U_j\|_{3,\infty}\asymp a(Nr^3)^{1/3}\asymp m,
\]
and
\[
 \|\nabla U_j\|_2^2\asymp Na^2r\asymp r^{-2}.
\]
Replication of the regular band gives total volume
\(Nr^3\asymp R^3\), fixed energy, and perimeter
\[
 Nr^2\asymp R^{11/6}j^{2/3}.
\]

For \(\rho\leq r\), boundedness of \(\phi\) gives
\[
 \int_{B_\rho(x)}|U_j|^2\,dx
 \lesssim a^2\rho^3\leq C\rho^\beta.
\]
At \(\rho=r\), this uses
\[
 a^2r^3\asymp N^{-1}
 \asymp m^{-1}j^{-2}\leq Cr^\beta
\]
for \(\beta<1/9\).  For \(r\leq\rho\leq d\), at most a bounded number of
packets contributes, so the same endpoint estimate applies.  For
\(d\leq\rho\leq1\), quasi-uniformity gives at most
\(C(1+N\rho^3)\) packets and therefore energy
\[
 C(N^{-1}+\rho^3)\leq C\rho^\beta.
\]
\end{proof}

\begin{proposition}[Instantaneous clock-rate sharpness]
\label{prop:driver}
At \(\Lambda=m^{7/3}j^{2/3}\), there are smooth compactly supported
divergence-free cores \(w_j\), of radius \(\Lambda^{-1}\) and amplitude
\[
 b=m\Lambda=m^{10/3}j^{2/3},
\]
such that
\begin{equation}
 \label{eq:driver-norms}
 \|w_j\|_{3,\infty}\asymp m,
 \qquad
 \|w_j\|_2^2
 \asymp\frac{m^2}{\Lambda}
 =m^{-1/3}j^{-2/3}\longrightarrow0,
\end{equation}
while, for a fixed cutoff multiple \(C\),
\begin{equation}
 \label{eq:driver-rate}
 \|P_{\leq C\Lambda}
 \PP\nabla\cdot(w_j\otimes w_j)\|_2
 \asymp m^2\Lambda^{3/2}.
\end{equation}
The cores satisfy \eqref{eq:packet-morrey} for every
\(\beta\leq1/7\).
\end{proposition}

\begin{proof}
Choose a fixed smooth compactly supported divergence-free
\(\phi\) for which
\(\PP\nabla\cdot(\phi\otimes\phi)\not\equiv0\), and let
\[
 w_j(x):=b\phi(\Lambda x).
\]
The first two claims follow directly from scaling.  The complete
nonlinear acceleration has \(L^2\) scale
\[
 b^2\Lambda\,\Lambda^{-3/2}
 =b^2\Lambda^{-1/2}
 =m^2\Lambda^{3/2}.
\]
Choose one fixed \(C\) large enough that the low-pass projection retains
a fixed fraction of the nonzero base-profile acceleration; scaling gives
\eqref{eq:driver-rate}.  The core energy at its own radius is
\[
 \frac{m^2}{\Lambda}
 =m^{-1/3}j^{-2/3}
 \leq C\Lambda^{-\beta}
\]
when \(\beta\leq1/7\).  The smaller- and larger-radius Morrey estimates
then follow as in Proposition~\ref{prop:packets}.
\end{proof}

\begin{remark}
Proposition~\ref{prop:packets} saturates the enstrophy power forced by
Corollary~\ref{cor:endpoint}.  Proposition~\ref{prop:driver} saturates
the \(m^2\Lambda^{3/2}\) instantaneous rate with vanishing energy.
Neither is a same-trajectory counterexample, and neither realises the
signed increment \eqref{eq:signed-high}.
\end{remark}

\section{Summability and residence audit}

The endpoint floor is not an integrated enstrophy floor.  Multiplication
by the entire fixed-block record gap gives only the formal scale
\begin{equation}
 \label{eq:enstrophy-weight}
 \Delta_j^{(L)}\Lambda_j^2
 \asymp R_j^{5/12}j^{1/3},
\end{equation}
which is summable.  Since
\[
 \int_{A_j}|u(t_j)|^3\,dx
 \asymp m_j^3
 \asymp R_j^{-3/2},
\]
the layer cubic and \(r_j\)-cutoff scales are
\begin{equation}
 \label{eq:cubic-weight}
 \Delta_j^{(L)}
 \int_{A_j}|u(t_j)|^3\,dx
 \asymp\frac{R_j^{5/4}}j
\end{equation}
and
\begin{equation}
 \label{eq:flux-weight}
 \Delta_j^{(L)}r_j^{-1}
 \int_{A_j}|u(t_j)|^3\,dx
 \asymp R_j^{1/12}j^{-1/3}.
\end{equation}
The viscous cutoff scale agrees with
\eqref{eq:enstrophy-weight}.  These are dimensional exponent audits,
not signed event lower bounds.  Every displayed sequence is
geometrically summable.

The physical turnover scale of the selected carrier is
\[
 \tau_j\asymp m_j^{-5}.
\]
In carrier coordinates, the forced ultraviolet length is
\begin{equation}
 \label{eq:ell-clock}
 \ell_j^{\mathrm{clk}}
 :=\frac{r_j}{R_j}
 \asymp m_j^{-1/3}j^{-2/3}.
\end{equation}
A fixed high-pass floor gives a fixed heat-filter subgrid-energy floor
at a comparable scale.  Thus a canonical first-crossing scale is no
larger than \(\ell_j^{\mathrm{clk}}\), up to constants and a fixed
reindexing if the earlier endpoint is selected.

If that first-crossing scale is comparable to
\(\ell_j^{\mathrm{clk}}\), the previously established record-residence
estimate has physical duration
\begin{equation}
 \label{eq:residence}
 \tau_j
 \frac{\ell_j^{\mathrm{clk}}}
 {\log(e+C/\ell_j^{\mathrm{clk}})}
 \asymp
 m_j^{-16/3}j^{-5/3}.
\end{equation}
The ratio of \eqref{eq:residence} to the record gap is
\begin{equation}
 \label{eq:residence-ratio}
 \asymp m_j^{1/6}j^{-2/3}\longrightarrow\infty.
\end{equation}
The ultraviolet energy would then be prestored across the block.  This
does not make its signed correlation fresh: phase or spatial alignment
may be recycled while the energy persists.  If the canonical scale is
far smaller than \eqref{eq:ell-clock}, even
\eqref{eq:residence-ratio} is unavailable.

\section{Findings, conjecture, and remaining obligations}

\subsection*{Robust findings, subject to external review}

\begin{enumerate}
\item The critical weak-\(L^3\) ceiling improves the global low-pass
nonlinear rate from the energy-only \(\Lambda^{5/2}\) power to
\(m_j^2\Lambda^{3/2}\).
\item A future fixed-energy amplitude layer is nearly empty at a fixed
earlier record, producing an order-one signed increment.
\item Frequencies below
\(\Lambda_j\asymp R_j^{-7/6}j^{2/3}\) cannot carry that increment.
\item One record endpoint has fixed ultraviolet \(L^2\) energy and
enstrophy at least \(R_j^{-7/3}j^{4/3}\).
\item Every Leray--Hopf continuation of this conditional branch loses a
fixed amount of kinetic energy at the first singular time.
\item The static packet and driver families show sharpness of the forced
spatial power and instantaneous clock rate under the inputs used.
\item All present exact \(q=4\) record-gap weights remain summable.
\end{enumerate}

\subsection*{Things still to prove}

\begin{enumerate}
\item Derive energy equality, or strong left \(L^2\) continuity, at
\(T^*\) from the exact \(q=4\) schedule; alternatively contradict its
failure using another Navier--Stokes law.
\item Exclude repeated phase or spatial recycling of a prestored
ultraviolet reservoir by a nonsummable signed, fresh, or orthogonal
charge.
\item Control canonical subgrid scales far below
\(\ell_j^{\mathrm{clk}}\).
\item Control the divergent-normalised-energy branch.
\item Treat Type-II schedules at or below the temporal five-power
barrier.
\item Prove one complete Clay alternative for arbitrary admissible data.
\end{enumerate}

\begin{conjecture}[No anomalous-loss exact \(q=4\) branch]
No Leray--Hopf continuation of a smooth finite-energy solution through
its first singular time can simultaneously have the exact \(q=4\)
first-record schedule and a fixed-energy near-extremising weak-\(L^3\)
layer.
\end{conjecture}

Corollary~\ref{cor:defect} proves the conjecture conditional on energy
equality at the first singular time.  That condition is not established
here.  No Clay alternative is proved.

\begin{thebibliography}{9}

\bibitem{CL20}
A.~Cheskidov and X.~Luo,
\emph{Energy Equality for the Navier--Stokes Equations in Weak-in-Time
Onsager Spaces},
Nonlinearity \textbf{33} (2020), 1388--1403.
\href{https://doi.org/10.1088/1361-6544/ab60d3}
{doi:10.1088/1361-6544/ab60d3};
arXiv:1802.05785v2.

\end{thebibliography}

\end{document}
